\subsection{Inside PyTorch: how the actor is built and how its weights change}
\label{sec:pytorch_actor}
This subsection opens the box: what the actor is in code, what PyTorch
computes when it is trained, and \emph{when} its weights change as new
simulation results arrive.

\paragraph{The network is a function with adjustable coefficients.}
In \code{train\_ppo.py} the actor is three linear layers with $\tanh$
between them:
\begin{lstlisting}[language=Python]
self.actor = nn.Sequential(
    nn.Linear(obs_dim, 64), nn.Tanh(),     # W1: 64 x obs_dim, b1: 64
    nn.Linear(64, 64),      nn.Tanh(),     # W2: 64 x 64,      b2: 64
    nn.Linear(64, act_dim))                # W3: act_dim x 64, b3: act_dim
self.log_std = nn.Parameter(torch.full((act_dim,), -0.5))
\end{lstlisting}
Written out, the actor is the formula
\begin{equation}
\mu_\theta(o)=W_3\,\tanh\!\big(W_2\,\tanh(W_1 o+b_1)+b_2\big)+b_3 ,
\label{eq:mlp}
\end{equation}
and $\theta$ is the list of all entries of $W_1,b_1,W_2,b_2,W_3,b_3$ plus
$\log\sigma$. For the single reservoir ($o$ has 7 entries, one action)
that is $7\cdot64+64+64\cdot64+64+64+1=4737$ coefficients plus $\log\sigma$.
For the cascade (16 observations, 2 actions) it is 5378 plus two
$\log\sigma$. The critic is a separate network of the same shape ending in
one output, $V_\phi(o)$. \code{nn.Linear} starts each weight as a small
random number (uniform in $\pm1/\sqrt{\text{inputs}}$). The untrained
actor is therefore a random, nearly flat function of the state, and its
releases are near the middle of the range.

It helps to think of Eq.~\eqref{eq:mlp} as an ordinary nonlinear
regression model with 4737 parameters. Training it is a calibration, but
with an unusual objective. We never have a ``measured'' correct action to
fit to. The objective is built from the simulated rewards instead.

\paragraph{From the network to a probability.}
\code{model.dist(obs)} returns the Gaussian
$\pi_\theta(a\mid o)=\mathcal N(\mu_\theta(o),\sigma^2)$, and
\code{log\_prob} evaluates, for a recorded action $a$,
\begin{equation}
\log\pi_\theta(a\mid o)=-\frac{(a-\mu_\theta(o))^2}{2\sigma^2}-\log\sigma-\tfrac12\log 2\pi
\label{eq:logp}
\end{equation}
(summed over the action entries when there are two). Everything PPO does
to the actor goes through this one number.

With one action (the single reservoir) this is a univariate Gaussian. With
two actions (the cascade, or a release plus a valve setting) it is a
bivariate Gaussian with a \emph{diagonal} covariance,
$\mathrm{diag}(\sigma_1^2,\sigma_2^2)$: one mean and one spread per
action and no correlation term, which is the same as two independent
Gaussians. In code, \code{torch.distributions.Normal(mu, log\_std.exp())}
acts element by element and \code{log\_prob(a).sum(-1)} adds the
per-action terms of Eq.~\eqref{eq:logp}. The two \emph{means} are still
coupled, because both come out of the same network from the same
observation. Only the exploration noise added to them is independent.

\paragraph{Where the spread $\sigma$ comes from.}
$\sigma$ is not computed from data. It is one more trainable parameter,
updated by the same gradient steps as the weights. The network stores
$\log\sigma$, one entry per action
(\code{self.log\_std = nn.Parameter(torch.full((act\_dim,), -0.5))}), and
uses $\sigma=e^{\log\sigma}$. Storing the logarithm keeps $\sigma$
positive whatever value the optimizer gives it, and makes its steps
relative: a step of 0.1 in $\log\sigma$ changes $\sigma$ by about 10\%.
$\sigma$ does not depend on the state. Exploration is equally wide whether
the reservoir is full or empty. A state-dependent or correlated spread is
possible, but this simpler form is the usual PPO default and was enough
here.

\emph{Starting value.} Plain PPO starts at $\log\sigma=-0.5$, i.e.\
$\sigma\approx0.61$ on the action range $[-1,1]$: wide exploration. After
behavior cloning (Section~\ref{sec:ppo_extras}) the start is set by hand
(\code{--bc-log-std}). Cloning is a least-squares fit of $\mu_\theta$ and
carries no information about $\sigma$, and a policy that already
imitates a good rule should explore only around it. The pond used
$\log\sigma=-1.2$ ($\sigma\approx0.30$).

\emph{How it is learned.} Because $\log\sigma$ appears in
Eq.~\eqref{eq:logp}, \code{loss.backward()} gives it a gradient like any
weight, and Adam updates it. For one action entry,
\begin{equation}
\frac{\partial\log\pi_\theta(a\mid o)}{\partial\log\sigma}
  =\frac{(a-\mu_\theta(o))^2}{\sigma^2}-1 ,
\label{eq:dlogsigma}
\end{equation}
which is positive for an action farther than one $\sigma$ from the mean and
negative for one closer. The loss weights each record by its advantage
$\hat A_t$. If the actions that turned out better than expected
($\hat A>0$) lie far from $\mu$, $\sigma$ grows (``exploring wider paid
off''). If they lie close to $\mu$, or the far-out actions were the bad
ones ($\hat A<0$), $\sigma$ shrinks (``stay near the mean''). $\sigma$
therefore settles where the advantage-weighted spread of actions says
further exploration no longer pays. We use no entropy bonus
(\code{--ent-coef 0}), so nothing else keeps $\sigma$ large, and in practice it shrinks
as the policy improves. In the cascade run, the \code{std} column of the
training log ($e^{\log\sigma}$) fell from 0.595 to 0.184 over 245 updates.

\emph{At evaluation} $\sigma$ is not used at all: the action is the mean,
$a=\mu_\theta(o)$.

\paragraph{One training step, line by line.}
The update loop in \code{train\_ppo.py} is:
\begin{lstlisting}[language=Python]
d     = model.dist(b_obs[idx])                 # forward pass: mu, sigma for 256 records
logp  = d.log_prob(b_act[idx]).sum(-1)         # Eq. (logp) for the recorded actions
ratio = (logp - b_logp[idx]).exp()             # rho = pi_theta / pi_old
pg_loss = -torch.min(ratio * b_adv[idx],
                     ratio.clamp(0.8, 1.2) * b_adv[idx]).mean()   # -L_clip
v_loss  = 0.5 * (model.value(b_obs[idx]) - b_ret[idx]).pow(2).mean()
loss = pg_loss + 0.5 * v_loss
opt.zero_grad()                                # forget the previous gradient
loss.backward()                                # dloss/dparameter for every parameter
nn.utils.clip_grad_norm_(model.parameters(), 0.5)
opt.step()                                     # Adam changes the parameters
\end{lstlisting}
\begin{enumerate}\itemsep2pt
\item \emph{Forward pass.} The 256 observations of the minibatch go
  through Eq.~\eqref{eq:mlp} at once (as a $256\times\text{obs\_dim}$
  matrix). PyTorch records every operation as it computes, building a
  \emph{computation graph} from the parameters to \code{loss}.
\item \emph{Inputs that are just data.} \code{b\_obs}, \code{b\_act},
  \code{b\_logp} (the log-probability under the policy that collected the
  data, $\pi_{\theta_{old}}$), \code{b\_adv} ($\hat A_t$) and \code{b\_ret}
  ($\hat R_t$) were computed before the update and are constants here. In
  particular the advantages come from the critic \emph{as it was} when the
  rollout ended. They are not recomputed while the actor changes.
\item \emph{Backward pass.} \code{loss.backward()} applies the chain rule
  through the recorded graph, from \code{loss} back to every parameter,
  and stores $\partial L/\partial\theta_i$ in each parameter's \code{.grad}.
  This is \emph{back-propagation}. It is the same thing as an adjoint
  sensitivity calculation, and it gives all 4737 derivatives for about the
  cost of one extra forward pass, instead of 4737 finite-difference runs.
\item \emph{Gradient clipping.} If the length of the whole gradient
  vector exceeds 0.5 it is scaled down to 0.5. Its direction is kept and
  its size is capped, a guard against one unusual minibatch throwing the
  weights far.
\item \emph{Adam step.} \code{opt.step()} changes each parameter (next
  paragraph). \code{zero\_grad()} is needed because PyTorch \emph{adds}
  new gradients to \code{.grad} rather than overwriting them.
\end{enumerate}

\paragraph{What the gradient says, for one record.}
Take one record with advantage $\hat A$ and suppose clipping is not active.
Its contribution to the actor loss is $-\rho\hat A$, and with
Eq.~\eqref{eq:logp} the chain rule gives
\begin{equation}
\frac{\partial(-\rho\hat A)}{\partial\theta}
= -\rho\,\hat A\,\frac{a-\mu_\theta(o)}{\sigma^2}\;\frac{\partial\mu_\theta(o)}{\partial\theta}.
\label{eq:actor_grad}
\end{equation}
Gradient descent moves $\theta$ \emph{against} this. If $\hat A>0$ (the
recorded action turned out better than the critic expected), the step
changes the weights so that $\mu_\theta(o)$ moves \emph{toward} the action
$a$ that was actually taken. If $\hat A<0$, it moves away from it. The
factor $\partial\mu_\theta/\partial\theta$ is what back-propagation
supplies. It says how each weight would have to change to move the
network's output at this particular state $o$. Because the weights are
shared by all states, the change also shifts the release at
\emph{similar} states. This generalization is what lets a few thousand
coefficients represent a rule over a continuous state space. Once
$\rho$ leaves $[0.8,1.2]$ in the direction the advantage favors, the
clipped term is the active one, its gradient is zero, and that record stops
pulling. The minibatch loss averages Eq.~\eqref{eq:actor_grad} over 256
records, so the step follows the consensus of many experiences, not any
single one. The spread $\log\sigma$ gets its own gradient from the same
loss. It shrinks when the advantage-weighted actions lie close to $\mu$,
which is how exploration narrows by itself as training proceeds.

\paragraph{The Adam update.}
Adam keeps two running averages for every parameter $\theta_i$: of its
gradient $g_i$ and of $g_i^2$. With $\beta_1=0.9$, $\beta_2=0.999$, step
count $k$ and learning rate $\alpha=3\times10^{-4}$:
\begin{align}
m_i &\leftarrow \beta_1 m_i+(1-\beta_1)g_i, &
v_i &\leftarrow \beta_2 v_i+(1-\beta_2)g_i^2,\nonumber\\
\theta_i &\leftarrow \theta_i-\alpha\,
  \frac{m_i/(1-\beta_1^k)}{\sqrt{v_i/(1-\beta_2^k)}+10^{-8}} .
\label{eq:adam}
\end{align}
$m_i$ is a smoothed gradient (momentum) and $\sqrt{v_i}$ its typical size,
so each parameter moves by roughly $\alpha$ per step in the direction its
recent gradients agree on, whatever its scale. A weight whose gradient
keeps flipping sign moves little. The divisions by $1-\beta^k$ correct the
start-up bias of the averages, which begin at zero. One optimizer (one set
of $m,v$) serves actor, critic and $\log\sigma$ together, and it persists
across updates, so momentum carries over from one batch of data to the
next.

\paragraph{When the weights change as new data arrive.}
The weights are \emph{not} updated after every simulated day. Training
alternates between two phases:
\begin{enumerate}\itemsep2pt
\item \emph{Collect, weights frozen.} The current weights are copied to
  the 16 simulator workers. Each runs 256 days with them, drawing random
  actions from $\mathcal N(\mu_\theta(o),\sigma^2)$ and recording $o$, $a$,
  $\log\pi_{\theta_{old}}(a\mid o)$, $V_\phi(o)$ and $r$. Nothing learns
  during this phase. The 4{,}096 new records are then sent back.
\item \emph{Learn, data frozen.} Advantages and returns are computed for the
  4{,}096 records (Section~\ref{sec:ppo_data}). The records are then
  reused for 10 epochs. Each epoch shuffles them into 16 minibatches of 256
  and takes one Adam step per minibatch, so there are 160 weight updates
  per batch of data. In the first minibatch $\theta=\theta_{old}$ and every
  $\rho=1$. As the steps accumulate, $\rho$ drifts away from 1 and the
  clipping starts to switch records off.
\end{enumerate}
After the 160 steps the records are \emph{discarded}, because they were
generated by a policy that no longer exists, and the new weights go back
to the workers for the next batch. Over the single-reservoir run this
cycle repeats 100 times: 409{,}600 simulated days and 16{,}000 Adam steps.
Each new batch of experience therefore enters the actor in three ways: as
the actions $a$ (where to move $\mu$), as the advantages $\hat A$ (which
way and how strongly), and, through the critic that is fitted on the same
records, as better advantages in the next batch.

\paragraph{Two related procedures use the same machinery.}
\emph{Behavior cloning} (Section~\ref{sec:ppo_extras}) trains the same
actor with the same Adam optimizer, but its loss is ordinary least
squares, $\frac1n\sum(\mu_\theta(o)-a^{\mathrm{rule}})^2$, with the rule's
actions as targets. That is regression in the plainest sense. At
\emph{evaluation} no gradients are computed (\code{torch.no\_grad()}):
the trained actor is just Eq.~\eqref{eq:mlp} evaluated with fixed
coefficients, and its output is the release decision.
